\documentclass[a4paper, serif, 10pt]{moderncv}

%% moderncv
% style and color
\moderncvstyle{banking}
\moderncvcolor{black}

% renew moderncv command to adapt for CJK colon
\renewcommand*{\cvitem}[3][.25em]{%
  \ifstrempty{#2}{}{\hintstyle{#2}：}{#3}%
  \par\addvspace{#1}}

\renewcommand*{\cvitemwithcomment}[4][.25em]{%
  \savebox{\cvitemwithcommentmainbox}{\ifstrempty{#2}{}{\hintstyle{#2}：}#3}%
  \setlength{\cvitemwithcommentmainlength}{\widthof{\usebox{\cvitemwithcommentmainbox}}}%
  \setlength{\cvitemwithcommentcommentlength}{\maincolumnwidth-\separatorcolumnwidth-\cvitemwithcommentmainlength}%
  \begin{minipage}[t]{\cvitemwithcommentmainlength}\usebox{\cvitemwithcommentmainbox}\end{minipage}%
  \hfill% fill of \separatorcolumnwidth
  \begin{minipage}[t]{\cvitemwithcommentcommentlength}\raggedleft\small\itshape#4\end{minipage}%
  \par\addvspace{#1}}

%% page layout/margins
\usepackage[top=2.5cm, bottom=2.5cm, left=1.5cm, right=1.5cm]{geometry}

\nopagenumbers{}

%% Babel config for English language
\usepackage[english]{babel}

%% fontspec
\usepackage{fontspec}

\IfFontExistsTF{Linux Libertine}{
  \setmainfont[Ligatures={TeX, Common}, Numbers=Lining, ItalicFont=Linux Libertine]{Linux Libertine}
}{}
\IfFontExistsTF{Linux Libertine O}{
  \setmainfont[Ligatures={TeX, Common}, Numbers=Lining, ItalicFont=Linux Libertine O]{Linux Libertine O}
}{}

%% CTeX
% CJK support, used to show CJK characters in the resume
%
% - fontset=none: disable builtin fontset but instead set the CJK font manually
% - heading=false: disable ctex heading
% - punct=kaiming: use kaiming punctuations styles for CJK
% - scheme=plain: use plain scheme, do not override \normalsize font size
% - space=auto: space settings for CJK characters
%
% ref:
% - http://ctan.mirrorcatalogs.com/language/chinese/ctex/ctex.pdf
\usepackage[UTF8, heading=false, punct=kaiming, scheme=plain, space=auto]{ctex}

\IfFontExistsTF{Noto Serif CJK SC}{
  \setCJKmainfont{Noto Serif CJK SC}
}{}
\IfFontExistsTF{Noto Sans CJK SC}{
  \setCJKsansfont{Noto Sans CJK SC}
}{}

%% line spacing
\usepackage{setspace}
\setstretch{1.125}

%% URL styling - use normal text instead of monospace
\urlstyle{same}

% Auto-underline all links
% Defer until \begin{document} so hyperref is already loaded
\AtBeginDocument{%
  \let\oldhref\href
  \renewcommand{\href}[2]{\oldhref{#1}{\underline{#2}}}%
}

%% Patch \httplink and \httpslink to handle full URLs
% The original moderncv macros always prepend http:// or https:// to URLs.
% This causes issues when the URL already has a protocol (e.g., https://example.com)
% resulting in malformed URLs like https://https://example.com.
% This patch checks if the URL already contains "://" and uses it directly if so.
% Note: We use \str_set:Nx (x-expansion) to expand macros like \@homepage first.
\ExplSyntaxOn
\str_new:N \g_yamlresume_url_str

\RenewDocumentCommand{\httpslink}{O{}m}{
  \str_set:Nx \g_yamlresume_url_str {#2}
  \str_if_in:NnTF \g_yamlresume_url_str {://}
    { \url{#2} }
    { \url{https://#2} }
}

\RenewDocumentCommand{\httplink}{O{}m}{
  \str_set:Nx \g_yamlresume_url_str {#2}
  \str_if_in:NnTF \g_yamlresume_url_str {://}
    { \url{#2} }
    { \url{http://#2} }
}
\ExplSyntaxOff

\name{佐藤 健太}{}
\title{シニアDevOpsエンジニア | クラウドインフラストラクチャスペシャリスト}
\phone[mobile]{+81 90-1234-5678}
\email{kenta.sato@example.jp}
\homepage{https://kenta-sato.dev/}

\address{日本、東京都、渋谷区、東京都渋谷区神南1-2-3、150-0041}{}{}

\extrainfo{{\small \faGithub}\ \href{https://github.com/kenta-sato-devops}{@kenta-sato-devops} {} {} {} • {} {} {} 
{\small \faLinkedin}\ \href{https://www.linkedin.com/in/kenta-sato-devops}{@kenta-sato-devops} {} {} {} • {} {} {} 
{\small \faTwitter}\ \href{https://twitter.com/kenta\_devops}{@kenta\_devops}}

\begin{document}

\maketitle

\section{基本情報}

\cvline{}{\begin{itemize}
\item クラウドインフラとDevOpsの設計・運用で8年以上の経験を持つエンジニア
\item AWS、Kubernetes、Terraformを活用したスケーラブルなシステム構築が得意
\item CI/CD、監視、セキュリティ、SREプラクティスを通じて開発生産性と信頼性を向上
\item チームリーダーとして後進の指導と組織横断的な自動化推進を経験
\item 海外チームとの連携や技術コミュニティでの登壇・執筆活動も積極的
\end{itemize}}
\section{学歴}

\cventry{2012年4月～2016年3月}
        {学士、情報工学、成績：3.7/4.0}
        {東京工業大学}
        {\url{https://www.titech.ac.jp/}}
        {}
        {\begin{itemize}
\item 情報工学を専攻し、システム開発とネットワークの基礎を習得
\item 分散システムとクラウドコンピューティングの研究に従事
\item 卒業研究ではKubernetesを用いたマイクロサービス基盤の性能評価を実施
\end{itemize}
\textbf{コース}：オペレーティングシステム、コンピュータネットワーク、データベースシステム、分散システム、クラウドコンピューティング、情報セキュリティ、アルゴリズムとデータ構造}
\section{職歴}

\cventry{2021年4月～現在}
        {シニアDevOpsエンジニア}
        {株式会社クラウドネクサス}
        {\url{https://www.cloudnexus.co.jp/}}
        {}
        {\begin{itemize}
\item AWSとKubernetesを中心としたクラウドインフラの設計・運用を主導
\item TerraformによるIaCを推進し、インフラ構築時間を70\%削減
\item GitHub ActionsとArgo CDを用いたCI/CDパイプラインを構築
\item PrometheusとGrafanaによる監視体制を整備し、MTTRを40\%改善
\item SREプラクティスを導入し、サービス信頼性と開発者体験を向上
\end{itemize}
\textbf{キーワード}：AWS、Kubernetes、Terraform、CI/CD、SRE、自動化、可観測性、インシデント対応}

\cventry{2018年4月～2021年3月}
        {DevOpsエンジニア}
        {株式会社デジタルフォージ}
        {\url{https://www.digitalforge.co.jp/}}
        {}
        {\begin{itemize}
\item オンプレミス環境からAWSへの移行プロジェクトを担当
\item DockerとKubernetesを用いたコンテナ基盤の構築・運用
\item JenkinsからGitLab CIへのCI/CD移行を主導
\item Ansibleによる構成管理の自動化と運用手順の標準化
\item 開発チームと協力し、リリースサイクルを月次から週次へ短縮
\end{itemize}
\textbf{キーワード}：Docker、Kubernetes、AWS、GitLab CI、Ansible、マイグレーション}

\cventry{2016年4月～2018年3月}
        {システムエンジニア}
        {日本テクノソリューションズ株式会社}
        {\url{https://www.ntsolutions.co.jp/}}
        {}
        {\begin{itemize}
\item 金融系システムの運用・保守を担当
\item ShellスクリプトとPythonによる運用自動化ツールを開発
\item 監視ツールの導入とアラート設計を実施
\item 障害対応と原因分析を通じてシステム安定性に貢献
\end{itemize}
\textbf{キーワード}：運用保守、Python、Shell、監視、障害対応}
\section{言語}

\cvline{日本語}{ネイティブ・母国語レベル \hfill \textbf{キーワード}：日本語ネイティブ}
\cvline{英語}{完全な実務レベル \hfill \textbf{キーワード}：TOEIC 900、技術文書読解、海外チームとの連携}
\section{スキル}

\cvline{クラウドプラットフォーム}{エキスパート \hfill \textbf{キーワード}：AWS、Azure、GCP、EKS、Lambda、CloudFormation}
\cvline{コンテナとオーケストレーション}{エキスパート \hfill \textbf{キーワード}：Docker、Kubernetes、Helm、Istio、Argo CD}
\cvline{IaCと自動化}{上級 \hfill \textbf{キーワード}：Terraform、Ansible、Pulumi、Packer、Shell、Python}
\cvline{CI/CD}{上級 \hfill \textbf{キーワード}：Jenkins、GitHub Actions、GitLab CI、CircleCI、Argo Workflows}
\cvline{監視と可観測性}{上級 \hfill \textbf{キーワード}：Prometheus、Grafana、Datadog、ELK、OpenTelemetry}
\cvline{セキュリティとコンプライアンス}{中級 \hfill \textbf{キーワード}：IAM、Vault、Trivy、OPA、SOC2、ISO 27001}
\cvline{プログラミング}{上級 \hfill \textbf{キーワード}：Python、Go、Bash、Ruby、SQL}
\section{受賞・表彰}

\cventry{2022年12月}
        {株式会社クラウドネクサス}
        {社内イノベーションアワード}
        {}
        {}
        {\begin{itemize}
\item Kubernetes基盤の標準化と自動化により、開発者のデプロイ時間を大幅に短縮
\item 全社的なDevOps文化の推進に貢献
\end{itemize}}

\cventry{2019年6月}
        {Amazon Web Services}
        {AWS Summit Japan Hackathon 優秀賞}
        {}
        {}
        {\begin{itemize}
\item サーバーレスアーキテクチャを用いた災害情報共有アプリを開発
\item チームリーダーとして設計とデプロイを担当
\end{itemize}}
\section{資格・免状}

\cventry{2022年6月}
        {Amazon Web Services}
        {AWS Certified DevOps Engineer – Professional}
        {\url{https://aws.amazon.com/certification/certified-devops-engineer-professional/}}
        {}
        {}

\cventry{2021年9月}
        {Cloud Native Computing Foundation}
        {Certified Kubernetes Administrator (CKA)}
        {\url{https://www.cncf.io/certification/cka/}}
        {}
        {}

\cventry{2021年3月}
        {HashiCorp}
        {HashiCorp Certified: Terraform Associate}
        {\url{https://www.hashicorp.com/certification/terraform-associate}}
        {}
        {}
\section{出版・著作}

\cventry{2023年9月}
        {クラウドネイティブ時代のKubernetes運用}
        {技術評論社}
        {\url{https://gihyo.jp/book/2023/978-4-297-13612-3}}
        {}
        {\begin{itemize}
\item Kubernetesクラスタの運用ノウハウとベストプラクティスを解説
\item 監視、ログ、セキュリティ、CI/CDの実践的手法を紹介
\item 共著者として第5章「可観測性とSRE」を担当
\end{itemize}}
\section{推薦人}

\cventry{\emaillink[seiichi.tanaka@example.jp]{seiichi.tanaka@example.jp}}
        {前職の技術部長}
        {田中 誠一}
        {+81 80-2345-6789}
        {}
        {佐藤健太氏は、クラウドインフラと自動化に関する深い知識を持ち、チームをリードしながら信頼性の高いシステムを構築できる優秀なエンジニアです。}
\section{プロジェクト}

\cventry{2022年4月～2023年3月}
        {オンプレミスからEKSへの移行を主導し、内製プラットフォームを構築}
        {社内Kubernetesプラットフォーム移行プロジェクト}
        {\url{https://github.com/kenta-sato-devops/eks-platform}}
        {}
        {\begin{itemize}
\item マルチテナントEKSクラスタを設計し、20以上のマイクロサービスを移行
\item TerraformとArgo CDによるGitOpsを導入
\item 移行後の運用コストを30\%削減
\end{itemize}
\textbf{キーワード}：Kubernetes、EKS、Terraform、GitOps、マイクロサービス}

\cventry{2021年10月～2022年3月}
        {PrometheusとGrafanaを用いた全社統合監視システムの構築}
        {統合監視ダッシュボード}
        {\url{https://github.com/kenta-sato-devops/observability-dashboard}}
        {}
        {\begin{itemize}
\item 複数チームの監視データを一元化し、障害検知を迅速化
\item SLOダッシュボードを整備し、サービス品質の可視化を実現
\end{itemize}
\textbf{キーワード}：Prometheus、Grafana、SLO、可観測性}
\section{興味・関心}

\cvline{オープンソース}{Kubernetes、Terraform、CNCF}
\cvline{アウトドア}{登山、キャンプ、ランニング}
\cvline{技術コミュニティ}{登壇、執筆、勉強会運営}
\section{ボランティア活動}

\cventry{2020年1月～現在}
        {運営メンバー}
        {Cloud Native Community Japan}
        {\url{https://community.cncf.io/cloud-native-community-japan/}}
        {}
        {\begin{itemize}
\item Kubernetes勉強会の企画・運営を担当
\item 初心者向けハンズオン資料を作成し、登壇者として発表
\item コミュニティイベントを通じてDevOps人材の育成に貢献
\end{itemize}}
    
\end{document}